\documentclass[11pt]{article}
\usepackage[margin=1in]{geometry}
\usepackage{amsmath,amssymb,mathtools,url}
\setlength{\parindent}{0pt}
\setlength{\parskip}{5pt}
\setlength{\emergencystretch}{2em}
\DeclareMathOperator{\Tr}{Tr}
\begin{document}
\paragraph{Causal sea, complete pair response and the scalar rows
[X/S; V:X,W,D].}
This result uses the original sine-polar $H_*$ and the auxiliary
mean-state time prescription of \path{solve/strand_X.tex}.
It introduces no coupling, shift, independent kinetic term or
target exterior.  The bounded questions are whether the
instantaneous protected vacuum has a controlled finite-frequency
replacement, which clock and ruler sources that replacement
generates, and which exact state equation supplies the alternative
when the replacement fails.  Vacuum response is distinguished from
finite-density matter response throughout.  No collapse or thermal
spectrum is inferred from a vacuum susceptibility.
Exactly two claim-specific [X] groups are established by the
independent methods stated with the sources: the completed-pair
causal ruler response with its static cancellation, and the
positive adiabatic clock source with its energy normalization.
The other checked calculations have supporting [S] scope.

\paragraph{1. Exact finite-parent state and source system.}
Use a finite native parent with fixed collar data and positive
numerical $N,\psi$.  Fermionic Fock space is finite-dimensional.
Bosonic Fock space is untruncated; require the common finite-particle
core and the state moments needed by all displayed forms and
derivatives.  The isolated limit requires its own source,
state and boundary convergence.
Write the original energy expectation as
\begin{equation}
 {\cal H}[N,\psi,\rho]
 =W_2[N,\psi]+{\cal R}_a[N,\psi]
       +\Tr(\rho\,H_{\rm exc}^{\rm nat}[N,\psi]),
 \qquad
 W_2=-k\int\nabla\psi\cdot\nabla(N\psi).
 \tag{XS1}
\end{equation}
Here $k=k_a(1)=(2\pi G)^{-1}$ is the single matched coefficient.
On a finite parent the integral denotes the specified primary
functional and its boundary convention; ${\cal R}_a$ is the
difference from the actual polar vacuum functional, with its
boundary allocation.  In the continuum primary approximation
use the continuum integral and keep the nonlocal,
four-derivative and $k_\circ-k_a(1)$ contributions in the remainder,
as in W33--W37.  Omitting this remainder is a named approximation.
It does not change the excitation Hamiltonian.

The state is native.  Field derivatives below hold that state
fixed:
\begin{equation}
 \begin{aligned}
 i\dot\rho&=[H_{\rm exc}^{\rm nat},\rho],\\
 C_x&=\partial_{N_x}(W_2+{\cal R}_a)
               +\Tr(\rho\,c_x)=0,\qquad
 c_x=\partial_{N_x}H_{\rm exc}^{\rm nat}\geq0,\\
 R_x&=\partial_{\psi_x}(W_2+{\cal R}_a)
               +\Tr(\rho\,r_x)=0,\qquad
 r_x=\partial_{\psi_x}H_{\rm exc}^{\rm nat}.
 \end{aligned}\tag{XS2}
\end{equation}
Only interior fields are varied.  All collar-crossing
coefficients and source labels remain.  The scalar vacuum
term can be included in the state generator; it contributes
only a phase.  The complete light forms are W9--W14,
including the full $n-1$ middle exponential in each $n$ sector.
The ruler derivative includes the moving excitation filter,
the regulator and distance kernel, the four-energy denominator,
the exponential derivatives and the native-frame commutator.
W5--W15 specify these terms without a free-sector substitution.

Let $T_\psi$ take a fixed excitation representation to native
CAR and set ${\cal K}=T_\psi^\dagger\dot T_\psi$.
The exactly equivalent state equation is
\begin{equation}
 i\dot\rho_{\rm pol}
 =[\,H_{\rm exc}-i{\cal K},\,\rho_{\rm pol}\,],
 \qquad \rho_{\rm pol}=T_\psi^\dagger\rho T_\psi .
 \tag{XS3}
\end{equation}
For a packet with oriented polar pair $(E,V)$, its one-particle
connection is the X6 matrix with
$EK+KE=V^\dagger\dot{\cal A}-\dot{\cal A}^\dagger V$.
The second-quantized connection retains pair creation,
annihilation and within-band terms.  It is included once.
Evolving $H_{\rm exc}$ alone in this moving representation
does not implement (XS2).

On a differentiable regular solution with fixed boundary data,
the same-functional identity is
\begin{equation}
 \dot{\cal H}=\sum_{x\in I}
           (\dot N_x C_x+\dot\psi_xR_x)=0.
 \tag{XS4}
\end{equation}
The state contribution is the trace of a commutator.
The clock functional is independent of $N$.  With
$J_{xy}=\partial_{\psi_y}C_x|_\rho$, its necessary tangency
equation is
\[
 \sum_{y\in I}J_{xy}\dot\psi_y
       +i\Tr\{\rho[H_{\rm exc}^{\rm nat},c_x]\}=0 .
\]
Differentiating the ruler row also retains $\dot N$, explicit
field derivatives and $i\Tr(\rho[H_{\rm exc}^{\rm nat},r_x])$.
The mixed field derivative is the transpose of $J$ on a
twice-differentiable common chart.  These equations diagnose
preservation; they do not prove invertibility, positive lapse
or a global collapsing branch.  In the isolated limit,
charge conservation and regional exported energy require
the boundary hypotheses of X3--X4.

\paragraph{2. Protected sea and the actual completed pair inverse.}
At a fixed background, let $\Omega_\psi$ be the normalized
protected polar sea and put
$Q_\Omega=1-|\Omega_\psi\rangle\langle\Omega_\psi|$.
The excitation energy annihilates it for every positive lapse.
For a real ruler profile $f$, define
\begin{equation}
 \eta_f=Q_\Omega\,\partial_f\Omega_\psi,\qquad
 O_f=\partial_f H_{\rm exc}^{\rm nat}.
 \quad
 O_f\Omega_\psi=-H_2\eta_f,\qquad
 c_x\Omega_\psi=0.
 \tag{XS5}
\end{equation}
The sea derivative is a particle-hole pair, apart from its
removable phase.  The full excitation law preserves excitation
number at fixed fields.  Consequently the vacuum transition
in (XS5) lies in its complete two-excitation space even though
$O_f$ is not generally a free bilinear.  The identity follows
by differentiating $H_{\rm exc}^{\rm nat}\Omega_\psi=0$.
Omitting a derivative in W13 need not preserve it.

For the light packet, the actual restriction and inverse are
\begin{equation}
 H_{2,\rm light}
  =e^{S_2}{\cal B}_{N,2}e^{S_2},\qquad
 H_{2,\rm light}^{-1}
  =e^{-S_2}{\cal B}_{N,2}^{-1}e^{-S_2},\qquad
 {\cal B}_{N,2}=d\Gamma_2(\operatorname{diag}(B_N,B_N^T)).
 \tag{XS6}
\end{equation}
This is the exact $n=2$ consequence of W12, since $S_1=0$.
It retains the actual full pp/ph/hh seed and its four
$B_1$ denominator.  The unchanged packets use their actual
free pair energies.  If $b=\lambda_{\min}({\cal B}_N)>0$,
then $H_{2,\rm light}\geq2b e^{-2\|S_2\|}$.
This finite-parent bound can be very small; the protected
one-excitation gap is not the response gap.
Let $\Delta_2>0$ be a valid lower bound on every pair block
reached by $\eta_f$.  No gap uniform in all excitation numbers,
volume or a vanishing-lapse limit is asserted.

The finite polar sea has a locally flat Berry connection.
For its occupied frame $W_-=(I,-V)^T/\sqrt2$,
$i\langle\Omega,d\Omega\rangle=(i/2)\Tr(V^\dagger dV)$,
summed over the actual fermionic packets.  It is locally
an exact determinant phase; its exterior derivative is zero
by the trace of the Maurer--Cartan commutator.  Empty fixed-CCR
bosons contribute no tangent.  This local phase choice does
not construct an infinite-sea Fock implementer.

\paragraph{3. Exact causal linear response for a real profile.}
Take a small prescribed variation $\phi(T)f$ about a fixed
background and a retarded preparation in its protected sea.
The linear moving-frame pair amplitude satisfies
\begin{equation}
 i\dot\zeta=H_2\zeta-i\eta_f\dot\phi,\qquad
 \zeta(T)=e^{-iH_2(T-T_0)}\zeta(T_0)
       -\int_{T_0}^T e^{-iH_2(T-t)}
                              \eta_f\dot\phi(t)\,dt .
 \tag{XS7}
\end{equation}
The initial term is retained unless the specified preparation
removes it.  This equation is exact to first order about
the fixed background.  On a finite-amplitude moving background,
the pair-preserving connection and higher-excitation couplings
in (XS3) cannot be dropped while calling a pair equation exact.
The same linear solution fixes its leading excitation sources:
\begin{equation}
 \begin{aligned}
 e_x^{(2)}(T)&=\langle\zeta,c_x|_2\zeta\rangle\geq0,\qquad
 r_f^{(1)}(T)=-2\operatorname{Re}
                         \langle\zeta,H_2\eta_f\rangle,\\
 \sum_xN_xe_x^{(2)}&=\langle\zeta,H_2\zeta\rangle,\qquad
 \frac{d}{dT}\langle\zeta,H_2\zeta\rangle
                              =\dot\phi\,r_f^{(1)} .
 \end{aligned}\tag{XS7a}
\end{equation}
The last identity is at the fixed background of (XS7).
It follows directly from that equation and is the causal
energy-work check.  The clock formula follows because every
$c_x$ annihilates the instantaneous reference, so neither
its normalization correction nor an unknown second-order
state amplitude contributes at this order.  It sums all
native source labels, including any collar.
These are linear-response/quadratic-energy identities,
not a nonlinear pair closure.

Let $d\mu_f(\lambda)=d\langle\eta_f,
{\bf1}_{(-\infty,\lambda]}(H_2)\eta_f\rangle$.
Differentiating the zero excitation eigenvalue twice gives
$\langle\partial_f^2H_{\rm exc}^{\rm nat}\rangle
=2\langle\eta_f,H_2\eta_f\rangle$.
This direct term cancels the static Lehmann term.  Thus for
$\operatorname{Im}z>0$, the diagonal real-profile ruler response is
\begin{equation}
 \begin{aligned}
 K^R_{ff}(z)
  &=(W_2+{\cal R}_a)_{,ff}-z^2 Z_f(z),\\
 Z_f(z)
  &=2\langle\eta_f,H_2(H_2^2-z^2)^{-1}\eta_f\rangle
    =2\int_{\Delta_2}^{\infty}
             \frac{\lambda\,d\mu_f(\lambda)}{\lambda^2-z^2},\\
 Z_f(0)&=2\langle\eta_f,H_2^{-1}\eta_f\rangle\geq0,\\
 Z_f(0)&\leq Z_f(\omega)
 \leq\frac{Z_f(0)}{1-\omega^2/\Delta_2^2}
                       \quad(|\omega|<\Delta_2).
 \end{aligned}\tag{XS8}
\end{equation}
The zero-frequency subtraction is supplied by the same
Hamiltonian, not a selected counterterm.  The light inverse
is (XS6).  Whole-packet orientation signs in the static
vacuum trace do not reverse these excitation spectral weights.
The complete clock has its signed polar row, but the
excitation clock vacuum-transition vertex is zero.
Its vacuum state-dependent linear clock response is therefore
zero.  Causal sea clock excitation begins at quadratic order.

Equation (XS8) is a causal linearization of the state-plus-row
system, or equivalently an in-in response.  It is not the
ordinary Hessian of a real single-history nonlocal action.
Off-diagonal profile response can have odd imaginary terms;
the flat Berry curvature removes the linear-frequency
antisymmetric term, not every higher odd term.
Only the stated diagonal real-profile theorem is used.

\paragraph{4. Exact finite-$q$ free-packet insertion.}
On a uniform prescribed background $\psi=P>0$, $N=N_0>0$,
let
\[
 k_p=P^{-2}d_a(p),\qquad
 \epsilon_p=(m^2+|k_p|^2)^{1/2},\qquad
 V_p=(m+i\sigma\cdot k_p)/\epsilon_p .
\]
All momenta lie in the one full native Brillouin zone.
For an additive ruler Fourier component $\phi_{p-p'}$,
polar Sylvester differentiation gives
\begin{equation}
 (V^\dagger\delta V)_{pp'}
       =\phi_{p-p'}\,{\cal A}_m(p,p'),\qquad
 {\cal A}_m(p,p')
 =\frac{(m-3i\sigma\cdot k_p)V_{p'}
       -V_p^\dagger(m+3i\sigma\cdot k_{p'})}
             {P(\epsilon_p+\epsilon_{p'})}.
 \tag{XS9}
\end{equation}
Indeed $\delta A_{pp'}=-P^{-2}(m+3i\sigma\cdot k_{p'})
\phi_{p-p'}$ and $E_p=\epsilon_p/P$ in the Sylvester equation.
The sea pair amplitude is $-{\cal A}_m/2$.
For the opposite orientation it is conjugated by the
endpoint polar unitaries and changes sign; its squared
Hilbert--Schmidt norm is identical.

With the Fourier convention
$V_{\rm box}^{-1}\int|\phi|^2=\sum_q|\phi_q|^2$,
the response density of a real Fourier profile is
\begin{equation}
 \begin{aligned}
 Z_m(q,z)&=\frac12\int_{\rm BZ}\frac{d^3p}{(2\pi)^3}
 \frac{\Delta(p,q)\,
       \Tr_2[{\cal A}_m(p+q,p)^\dagger{\cal A}_m(p+q,p)]}
                 {\Delta(p,q)^2-z^2},\\
 \Delta(p,q)&=N_0(\epsilon_{p+q}+\epsilon_p)\geq2N_0m .
 \end{aligned}\tag{XS10}
\end{equation}
Momenta are taken modulo the native zone.  Exchange of the
two endpoints gives the same integrated weight for $q$ and
$-q$.  A cosine/sine standing-wave profile therefore has
the real diagonal spectral form (XS8).
No derivative expansion in $q$ has been made.
For a finite periodic parent replace the integral by its
actual normalized momentum sum.  The bulk integral is its
thermodynamic limit, not a claim about every small grid.
The full light formula at momentum $q$ is (XS8) with the
actual pair-fiber inverse (XS6), not (XS10) at a dressed mass.

At $q=0$,
${\cal A}_m(p,p)=-2im\sigma\cdot k_p/(P\epsilon_p^2)$,
so (XS10) recovers exactly
\begin{equation}
 Z_m(0,0)=\frac{2m^2}{N_0P^2}
 \int_{\rm BZ}\frac{d^3p}{(2\pi)^3}
              \frac{|k_p|^2}{(m^2+|k_p|^2)^{5/2}} .
 \tag{XS11}
\end{equation}
This is the additive $\delta\psi$ response; a
$\delta\log\psi$ perturbation multiplies it by $P^2$.
Equations (XS9)--(XS11) keep the sine law, spin trace,
all valleys and finite spacing explicitly.

\paragraph{5. A finite-spacing lower bound and the scale comparison.}
Put $\delta=amP^2$ and
\[
 {\cal L}(x)=\operatorname{arsinh}x
 -\frac{x}{\sqrt{1+x^2}}
 -\frac{x^3}{3(1+x^2)^{3/2}},\qquad x\geq0 .
\]
The elementary positive derivative is
${\cal L}'(x)=x^4/(1+x^2)^{5/2}$.
Changing to $u=ap$, folding the 8 sine cells, and then
using $v_i=\sin u_i$ gives
\begin{equation}
 \begin{aligned}
 Z_m(0,0)&=\frac{2m^2P^4}{N_0}
 \int_{[-\pi,\pi]^3}\frac{d^3u}{(2\pi)^3}
 \frac{\sum_i\sin^2u_i}
       {(\delta^2+\sum_i\sin^2u_i)^{5/2}}\\
 &\geq\frac{8m^2P^4}{\pi^2N_0}\,
                                  {\cal L}(\delta^{-1}) .
 \end{aligned}\tag{XS12}
\end{equation}
The inequality restricts the folded $v$ cube to its unit
ball and uses
$\prod_i(1-v_i^2)^{-1/2}\geq1$.
It is a thermodynamic full-zone bound at every positive
$a,m,P,N_0$, not an extra valley multiplier.

At fixed positive $m,P,N_0$, the same integral yields
\begin{equation}
 Z_m(0,0)=
 \frac{8m^2P^4}{\pi^2N_0}
       \log\frac1{amP^2}
       +O(m^2P^4/N_0)\qquad(a\to0).
 \tag{XS13}
\end{equation}
To derive the coefficient, isolate small neighborhoods of
the 8 sine zeros in local periodic coordinates.
There $\sin u=u+O(|u|^3)$, and the leading radial integral
is $\int r^4(\delta^2+r^2)^{-5/2}dr$.
The Jacobian and sine corrections have a bounded integral,
as does the complement of the neighborhoods.
Thus the coefficient is fixed by the native spin count and
zone, without a continuum cutoff fit.

At fixed lapse the primary ruler Hessian has magnitude
$2kN_0q^2$.  Define its real-profile comparison
\[
 {\mathfrak R}(q,\omega)
       =\frac{\omega^2Z_{\rm vac}(q,\omega)}
                        {2kN_0q^2}.
\]
First take $q\to0$ at fixed positive inputs with
$\omega=vN_0|q|/P^2$ and fixed real $v$.
This is subgap for sufficiently small $q$.
The parameter $v$ specifies a comparison rate, not an imposed
material or surface velocity.  For a statement about the total
response, require existence of its homogeneous pair-response
limit and a positive reached completed-light gap uniformly
near $q=0$ in that thermodynamic or box limit, at this fixed
spacing.  A positive gap on each finite parent alone does
not imply this premise.  Without those inputs the following inequality
is a result for the unchanged heavy contribution alone.
The unchanged heavy packets alone then give
\begin{equation}
 \liminf_{q\to0}{\mathfrak R}(q,\omega)
 \ \geq\
 \frac{4v^2}{\pi^2k}
       \sum_{s\ {\rm unchanged\ heavy}}m_s^2
                  {\cal L}((am_sP^2)^{-1}) .
 \tag{XS14}
\end{equation}
The interacting light vacuum adds a nonnegative diagonal
subgap weight, provided its actual pair threshold is respected.
For a fixed positive mass rectangle with
$k_a(1)\to k_\circ=3J/\pi^2$, the right-hand side grows
logarithmically as $a\to0$ when $v\ne0$.
No coefficient is refitted in taking this diagnostic limit.
At a fixed physical spacing (XS12),(XS14), or the full
integral (XS10), is the test; the asymptotic is not a
numerical evaluation there.

This proves that a small frequency relative to a packet gap
does not uniformly control the instantaneous approximation
at a light-cone-scale rate.  It does not prove a scalar
propagating mode, instability, collapse rate or frozen surface.
The clock constraint also has to hold.  In particular, at a
homogeneous vacuum background its linearization is
$-kPq^2\phi+\delta e_{\rm mat}=0$.
There is no independent vacuum conformal oscillation supplied
by the inertia alone.
The actual finite-$a$ spatial remainder and an inhomogeneous
response enclosure are separate from (XS14).

\paragraph{6. The adiabatic sources and a direct validity bound.}
In a smooth finite-parent chart let
$\eta_x=Q_\Omega\partial_{\psi_x}\Omega$ and
$\eta_{\dot\psi}=\sum_x\dot\psi_x\eta_x$.
The local subgap state elimination gives
\begin{equation}
 \begin{aligned}
 L_{\rm ad}&=\langle\eta_{\dot\psi},
                          H_2^{-1}\eta_{\dot\psi}\rangle
      =\frac12Z_{xy}\dot\psi_x\dot\psi_y,\qquad
 Z_{xy}=2\operatorname{Re}\langle\eta_x,H_2^{-1}\eta_y\rangle,\\
 \xi&=H_2^{-1}\eta_{\dot\psi},\qquad
 -\partial_{N_x}L_{\rm ad}
       =\langle\xi,c_x|_2\,\xi\rangle\geq0 .
 \end{aligned}\tag{XS15}
\end{equation}
The reference sea and its ruler tangents are independent
of $N$.  Differentiating the inverse, with
$\partial_{N_x}H_2=c_x|_2$, proves the last identity.
Independently, the normalized state proportional to
$\Omega+i\xi$ has excitation clock expectation
$\langle\xi,c_x\xi\rangle/(1+\|\xi\|^2)$.
It agrees through quadratic order in ruler velocity.
This is a same-law positive source, not a negative vacuum
subtraction or an added matter species.
The homogeneity of $H_2$ in the complete lapse vector gives,
with all native source labels including the collar retained,
\[
 \sum_xN_x(-\partial_{N_x}L_{\rm ad})
     =\langle\xi,H_2\xi\rangle=L_{\rm ad}.
\]
There is no factor of two in this weighted clock energy.
The nonlinear lapse dependence of $L_{\rm ad}$ (uniformly
proportional to $N_0^{-1}$) results from eliminating the
causal state.  The original Hamiltonian remains affine in
$N$ at fixed native state.

If a matter functional ${\cal F}_{\rm mat}$ has been separated
from the eliminated vacuum modes with its approximation
controlled, the effective field rows obtained from
$L_{\rm ad}-(W_2+{\cal R}_a+{\cal F}_{\rm mat})$ are,
at this derivative order and in the local Berry phase chart,
\begin{equation}
 \begin{aligned}
 0&=C_{{\rm vac},x}+e_{{\rm mat},x}
                         +\langle\xi,c_x|_2\xi\rangle,\\
 0&=R_{{\rm vac},x}+r_{{\rm mat},x}
       +\frac{d}{dT}(Z_{xy}\dot\psi_y)
       -\frac12(\partial_{\psi_x}Z_{yz})\dot\psi_y\dot\psi_z .
 \end{aligned}\tag{XS16}
\end{equation}
Repeated site labels are summed.  The time derivative
retains $\dot N\,\partial_N Z$ as well as ruler derivatives.
The coefficient derivative in the native chart is
\[
 \partial_f L_{\rm ad}
  =2\operatorname{Re}
       \langle\partial_f\eta_{\dot\psi},\xi\rangle
       -\langle\xi,(\partial_fH_2)\xi\rangle ,
\]
at fixed velocities.  Derivatives of the embedded pair
subspace and the native-frame connection are included;
in a transported frame the corresponding covariant
derivatives replace the ordinary ones.
The actual $\partial_fH_2$ differentiates both exponentials
in (XS6), the $B_N$ matrix and W9's full $S_2$ derivative.
The moving-filter term is neither omitted nor added twice.
There is no general sign for the ruler row in (XS16);
the free principal comparison X9 consequently does not
apply to this causal source by substitution.
In the locally parallel flat-Berry chart the excitation
ruler force in (XS16) can also be written
\[
 R_{{\rm ad},x}^{\rm exc}
   =\langle\xi,r_x\xi\rangle
                   +2\operatorname{Re}\langle\eta_x,\dot\xi\rangle .
\]
This follows by differentiating
$\partial_{\dot\psi_x}L_{\rm ad}
=2\operatorname{Re}\langle\eta_x,\xi\rangle$ and using
$\partial_{\psi_x}\eta_y=\partial_{\psi_y}\eta_x$ in this
chart.  It makes the acceleration contribution explicit.

For a check not based solely on derivative counting, fix
prescribed differentiable fields and choose the parallel sea
phase.  Remove the common scalar $W_0$ phase from the native
state equation.  With $\Phi_{\rm ad}=\Omega+i\xi$,
\begin{equation}
 (i\partial_T-H_{\rm exc}^{\rm nat})\Phi_{\rm ad}
       =-\dot\xi,\qquad
 \|\Psi(T)-\Phi_{\rm ad}(T)\|
 \leq \|\Psi(T_0)-\Phi_{\rm ad}(T_0)\|
                  +\int_{T_0}^T\|\dot\xi(t)\|dt .
 \tag{XS17}
\end{equation}
This Duhamel estimate uses the full native propagator.
$\Phi_{\rm ad}$ is deliberately unnormalized in the identity;
normalization changes it by at most $\|\xi\|^2/2$.
If the displayed error is $\varepsilon$, a bounded source $O$
has expectation error at most
$\|O\|\varepsilon(2\|\Phi_{\rm ad}\|+\varepsilon)$ before
the stated normalization correction.
For the ruler source this estimates the trial-state expectation
$\langle\Phi_{\rm ad},r_x\Phi_{\rm ad}\rangle
=\langle\xi,r_x\xi\rangle$ in the parallel chart.
It does not by itself estimate the whole force in (XS16).
A comparison with that force additionally budgets
$2|\operatorname{Re}\langle\eta_x,\dot\xi\rangle|$ and the
normalization error.  Alternatively its state approximation
must include the next orthogonal correction
$-H_{{\rm exc},\perp}^{-1}Q_\Omega\dot\xi$ and bound that
approximation's residual, with its phase and normalization.
No such higher-order residual estimate is asserted here.
The derivative $\dot\xi$ is the full native derivative,
including its components outside the instantaneous pair space.
The estimate is useful only with an actually small bound over
the time interval; it is not a uniform late-time adiabatic
claim.  Unbounded bosonic observables need their corresponding
moment/form bounds.  The empty reference bosons have no
contribution to (XS15).

Equations (XS15)--(XS17) are a vacuum elimination test, not
a formula for the inertia of an occupied dressed star.
Finite-density matter needs its own state response and W's
complete clock/ruler expectation; adding a vacuum formula
to a shared occupied sector does not establish that separation.
If the state and force bounds above are not small,
retain (XS2)--(XS3), rather than both adding $L_{\rm ad}$
and evolving the same sea degrees of freedom.

\paragraph{7. What pair production this calculation determines.}
For a real pulse $\phi(T)f$ with $\phi(T)=\epsilon F(T)$,
where $F$ is a fixed sufficiently smooth compact pulse,
start in the protected sea and define
$\widehat\phi(\lambda)=\int e^{i\lambda T}\phi(T)\,dT$.
First-order transition amplitudes from (XS7) give
\begin{equation}
 {\cal P}_2=\int \lambda^2|\widehat\phi(\lambda)|^2
                         d\mu_f(\lambda)+O(\epsilon^3),\qquad
 {\cal E}_2=\int \lambda^3|\widehat\phi(\lambda)|^2
                         d\mu_f(\lambda)+O(\epsilon^3).
 \tag{XS18}
\end{equation}
Here ${\cal P}_2$ is probability in the two-excitation space
at this order, not an outgoing particle flux.
The remainder is as $\epsilon\to0$ with fixed finite-parent
positive-field/domain bounds and fixed duration.
It is not uniform in duration, volume, a closing gap or an
infinite collapse history.
For a free packet its massive pair threshold is the actual
$\Delta(p,q)$ in (XS10).  For the light packet use its
completed spectral measure, which can include bound pair
states.  Outgoing flux further requires the channel/scattering
readout, geometry and preparation.  The exact sine kernel
permits a controlled low-$a$ expansion only after its occupied
momentum and frequency ranges are bounded.  No universal
thermal factor or order-$(\kappa a)^2$ outgoing coefficient
follows from (XS18) without the nonlinear history and matching.

\paragraph{Constructive collapse handoff and disposition.}
Use the same boundary clock $T$ and regular solution of
(XS2); evaluate matter and vacuum departure in that same
native state.  At finite coupling the clock is W14 and the
ruler source is W5--W15.  The primary radial field rows are
W35 with the stated remainder derivatives.  Solve the clock
row for $\psi$, the ruler row for $N$, and evolve the complete
native state, or equivalently (XS3).  The regularity, positive
lapse and isolated-limit controls are inputs to this
differential-algebraic evolution, not conclusions of its notation.

The bounded response test is (XS8)--(XS14) at the retained
model masses, spacing and couplings, with the actual completed
light pair inverse.  A free heavy lower contribution already
of order the spatial stiffness rejects neglect of causal
vacuum response at the stated rate.  The adiabatic alternative
requires both the true pair-gap condition and the state/source
error (XS17), and includes both rows of (XS16).  It cannot be
adopted solely because $|\omega|a$ is small.

Monitor the same-functional boundary charge and exported
energy; both exterior coefficients and the conditional
$\dot D=(B-A)/G$ identity with W's remainder/boundary defects;
the full constraint residuals and tangency equations;
positive lapse and field Jacobian; material-charge quantiles
and resolved widths; $N\psi$ and the native speed
$N\psi^{-2}$ as well as the time-dependent areal characteristics;
the native momentum window; sea departure, pair spectral
weight and outgoing readout; and the actual response/spatial
stiffness ratio.  A lapse zero, target lock or material speed
is not inserted into the equations.

Causal sea response can supply a missing ruler acceleration
and positive clock energy within the same $H_*$.  The exact
clock identity and the response bounds show what this
statement means; they do not establish that it supplies the
required nonlinear surface behavior.  If the causal regular
solution still exhibits resolved finite native light speed
at the charge-scale material surface, the next named target
is the complete matter source or this auxiliary mean-state
time prescription.  A law change would have to be declared
before its consequences.  No such change is made here.

\paragraph{Scope and SOURCES.}
[X/S; V:X,W,D.]  The result gives the finite-parent exact
state/row coupling, a causal real-profile response with the
actual light pair inverse, an exact finite-$q$ free-packet
formula, a heavy-packet lower bound and logarithmic comparison,
and the induced positive adiabatic clock source with its
ruler companion and residual test.  These are vacuum response
statements with specified domains.  They narrow the validity
of the instantaneous-vacuum approximation and prescribe its
same-law causal replacement.  They do not prove a nonlinear
dressed collapsing branch, a charge-radius locking theorem,
an isolated-limit dynamics or a thermal spectrum.
The completed-pair causal ruler kernel and its static
cancellation (XS5)--(XS8) have [X] scope from differentiation
of the protected family followed by the Lehmann response,
and independently the moving-sea Volterra equation with its
direct force/work identity.
The positive adiabatic clock source and weighted-energy
normalization (XS15) have [X] scope from lapse differentiation
of the inverse, and independently the normalized corrected-state
expectation with the lapse-affine energy sum.
The finite-$q$ Sylvester normalization, sine lower bound and
logarithm, full-source transport, Duhamel and force budgets,
pulse production and boundary/characteristic conditions are
checked supporting [S] under their stated premises.
The auxiliary time prescription, thermodynamic limits,
uniform completed-light subgap condition, in-medium response
and nonlinear collapse/production applications retain
[V:X,W,D].  There is no new [R] or [N] result.

The original completion and protected reference are
\path{solve/strand_A_lattice.tex}, (A3)--(A7),(A14)--(A18).
The exact finite-coupling derivative, denominator and middle
factor are \path{solve/strand_W.tex}, W5--W15; its conditional
continuum equivariance and primary/remainder distinction are
W26--W37, independently reviewed in
\path{solve/review_W.tex}.
The time prescription, boundary/characteristic conditions and
principal shell comparison are \path{solve/strand_X.tex} and
\path{solve/review_X.tex}.
The full pair-response and adiabatic-inverse input
\path{solve/strand_N_vacuum_collective.tex}, N146--N148,
is rederived here through (XS5)--(XS8),(XS15).
The actual derivations of this result are (XS1)--(XS18).
The single independent substantive review is
\path{solve/review_X_causal_sea.tex}; its source and domain
repairs are incorporated here.
No web-return assertion is used as a proof or numerical input.
\end{document}
